  \begin{definition} \textbf{Conjunto abierto.} \label{def:abierto}
       Sea $A  \subseteq  \Rn{n}$. Decimos que $A$ es \emph{abierto} si cumple que
   \[
        \forall \vect{x}_0 \in \exists \delta > 0 : B_{\delta}(\vect{x}_0)  \subseteq  A.  
%\footnote{\text{\textit{i.e.}: $A$  es abierto si todos sus puntos son interiores.}}
   \]
   \end{definition}
   
   \begin{propertie}  \label{prop:abierto_int}    Sea  $A \subseteq \Rn{n}$, entonces:
   \begin{enumerate}
    \item \( \interior{A} \subseteq A \).
    \item \(A\) es abierto \(\iff A=\interior{A}\).
    \end{enumerate}
    \end{propertie}
   
   \begin{definition} \textbf{Conjunto cerrado.} \label{def:cerrado}
    Sea $A  \subseteq  \Rn{n}$. Decimos que $A$ es \emph{cerrado} si su complemento $A^c = \Rn{n} \smallsetminus A$ es abierto.
   \end{definition}
   